% =============================================================================
%  Tiểu kết chương 4
% =============================================================================

\section*{Tiểu kết}
\addcontentsline{toc}{section}{Tiểu kết}

Chương này đã trình bày quá trình thực nghiệm và đánh giá hệ thống GraphRAG đề xuất thông qua hai hướng tiếp cận gồm benchmark tự động trên bộ dữ liệu CheoBench và thực nghiệm với người dùng thực tế. Kết quả benchmark trên 100 ca kiểm thử cho thấy GraphRAG đạt hiệu năng cao hơn so với RAG truyền thống, chủ yếu ở các độ đo liên quan đến chất lượng ngữ cảnh và độ tin cậy của câu trả lời. Dù hai hệ thống có kết quả tương đương ở các độ đo IR thuần như MAP và NDCG@10, GraphRAG thể hiện ưu thế hơn trong việc duy trì tính đầy đủ và liên kết của ngữ cảnh truy xuất, qua đó cải thiện chất lượng tạo sinh câu trả lời cuối cùng.\n\nCác thực nghiệm phân tích thành phần cũng cho thấy hiệu quả của kiến trúc đề xuất. Chiến lược truy xuất đa mức giúp nâng cao khả năng tổng hợp tri thức giữa các thực thể và cộng đồng liên kết, trong khi định dạng biểu diễn ngữ cảnh có cấu trúc, đặc biệt là JSON, hỗ trợ mô hình ngôn ngữ khai thác quan hệ tri thức hiệu quả hơn so với các biểu diễn tuần tự thuần văn bản. Đồng thời, việc lựa chọn các chiến lược tạo sinh linh hoạt cũng cho thấy hiệu quả cao hơn khi hệ thống xử lý đa dạng các loại câu hỏi. Những kết quả này cho thấy hiệu quả của GraphRAG không chỉ đến từ truy xuất tài liệu mà còn từ khả năng tổ chức và biểu diễn tri thức theo cấu trúc đồ thị.\n\nThực nghiệm với người dùng cũng xác nhận tính khả thi của hệ thống trong môi trường sử dụng thực tế. GraphRAG đạt tỷ lệ ưu tiên lựa chọn cao nhất ở cả hai hình thức đánh giá, đồng thời được người dùng đánh giá cao hơn về độ chính xác, tính đầy đủ và mức độ đáng tin cậy của câu trả lời. Kết quả này cho thấy kiến trúc GraphRAG phù hợp với các miền tri thức có cấu trúc quan hệ phức tạp như nghệ thuật Chèo, nơi việc tổng hợp và liên kết tri thức đóng vai trò quan trọng hơn truy xuất văn bản đơn thuần.\n\nBên cạnh các kết quả tích cực, hệ thống vẫn còn một số hạn chế như chi phí thời gian phản hồi tương đối cao do quá trình truy xuất và mở rộng đồ thị, đồng thời chưa có bước đánh giá chuyên sâu bởi các chuyên gia trong lĩnh vực nghệ thuật Chèo. Đây sẽ là những hướng cần tiếp tục hoàn thiện trong các nghiên cứu tiếp theo.
